\Needspace{6\baselineskip}
\section{Introduction}
\IEEEPARstart{P}{hysics}-based design automation must turn requirements into a configuration that can be compiled, simulated, and checked. In a data center, this includes the arrangement and orientation of racks and cooling equipment, its supply/return architecture, and operating parameters. Similar structure appears in aerodynamic and thermal component design: a discrete choice determines the design family, while continuous parameters determine its geometry or operating point. A language model offers a common interface for requirements and structured proposals, and post-training can adapt this interface to repeated design tasks.

The cost of physical feedback is a central constraint. A reduced model or coarse simulation can score many designs, while a detailed simulation may require substantial computation and geometry preparation. Evaluators also differ in the information they expose. A black-box solver returns a score; a diagnostic interface returns component losses; a residual interface checks a proposed field against governing equations; and a differentiable or adjoint interface returns sensitivities. This creates an opportunity to purchase and combine several kinds of feedback instead of treating every evaluator call as an interchangeable scalar reward.

The learning objective is nevertheless defined by a particular high-fidelity evaluator. Directly maximizing a cheap score can favor configurations whose apparent improvement comes from model discrepancy. High-fidelity checks after training reveal this discrepancy but do not themselves determine how earlier updates should have used cheap information. Conversely, evaluating every sampled design at high fidelity may leave too few updates to learn a reusable policy. The relevant question is how available physical information can reduce the cost of learning toward the specified target.

We propose \emph{information-stratified policy optimization} (ISPO), a post-training framework that jointly manages physical fidelity, information content, and update cost. It uses cheap physical signals to predict a target-consistent update, then corrects the prediction with a randomized subset of high-fidelity evaluations. Recording the sampling probabilities allows objective-aware selection without silently changing the training distribution. Continuous sensitivities enter through a numerical output head; discrete structure remains an explicit stochastic decision.

The paper develops three connected claims: physical information can improve the economy of target-directed post-training; objective-aware residuals can support reliable learning with imperfect approximations; and the resulting policy can reuse physical design knowledge across requirements. ISPO addresses these claims through information-aware control variates, residual-guided acquisition, and distinct update routes for structural and numerical decisions. Its conditions for target consistency are stated explicitly. Each claim has a corresponding experiment: information and cost ablations test learning economy, fidelity-gap studies test the role of residuals, and non-language-policy, search, and held-out design comparisons test the value of reusable post-training. The complete experimental program is retained as the evidence needed to establish these claims.

Data-center design is the motivating application. A controlled thermal benchmark provides inexpensive reference evaluations and tunable approximation errors; an aerodynamic benchmark supplies an additional problem with a mature adjoint interface. This organization makes the method testable before undertaking large layout experiments, while keeping the intended contribution in post-training for physical design automation.
